\begin{proof}[Proof of \cref{obj:thm:sharp-local-linear-constant-and-representers}]
\begin{enumerate}
\item Fix \(t\) and write \(n=n_t\), \(d=d_t\), and \(m=m_t\).  On the complete-block graph, every active block is a full neighborhood of size \(d\), and \(n=md\).  The canonical block-local weights
\[
w^0_{i,t}(Z)=g_d(Z_{b(i)})
\]
belong to \(\mathcal E_t^{\mathrm{loc,lin}}\).  The block program identities in \cref{obj:lem:block-energy-representer}, applied block by block, give value \(A_d\) per unit score and therefore \(d^2A_d\) per block in the block-extremal criterion.  Thus every admissible local-linear rule has worst-case risk at least
\[
\frac{B^2}{n^2}\sum_{b=1}^{m}d^2A_d
 = \frac{B^2A_d}{m},
\]
while the canonical rule attains this value.  Since \(n=md\),
\[
\frac{B^2A_d}{m}=\frac{B^2dA_d}{n}.
\]
This proves the finite minimax identity. % lean: locLinMinimaxRisk_exact

\item For a fixed \(w\in\mathcal E_t^{\mathrm{loc,lin}}\), the squared error decomposes across complete blocks.  Maximizing the coefficient signs and low-order monomials independently on each block gives
\[
\sup_{c\in\mathcal M_t(G_t)}
\mathbb E_Z[(\widehat\tau_{w,t}-\tau_{n_t})^2]
=
\frac{B^2}{n_t^2}\sum_{b=1}^{m_t}\Psi_{b,t}(w),
\]
where the finite maximum in \(\Psi_{b,t}\) is the same block-extreme functional after translating the sign coding to \(\varepsilon_i\in\{-1,1\}\) and writing block monomials as \(Z_{S_i}\). % lean: locLinWorstRisk_exact_blockExtremal

\item The per-block lower bound follows from the all-empty monomial candidate and the block score identity.  Let \(g\) be the common complete-block score on \(V_b\).  For every \(i\in V_b\), blockwise unbiasedness gives
\[
\mathbb E_Z[w_{i,t}(Z_b)g(Z_b)]=A_{d_t},
\qquad
\mathbb E_Z[g(Z_b)^2]=A_{d_t}.
\]
Hence
\[
0\le
\mathbb E_Z\!\left[
\left\{\sum_{i\in V_b}w_{i,t}(Z_b)-d_t g(Z_b)\right\}^2
\right]
=
\mathbb E_Z\!\left[
\left\{\sum_{i\in V_b}w_{i,t}(Z_b)\right\}^2
\right]
-d_t^2A_{d_t}.
\]
The all-empty choice in \(\Psi_{b,t}\) contains the first expectation on the right, so
\[
d_t^2A_{d_t}\le \Psi_{b,t}(w).
\]
% lean: completeBlock_blockExtremal_lower

\item Define
\[
Y_t:=
\frac{1}{m_t d_t^2 A_{d_t}}
\sum_{b=1}^{m_t}\{\Psi_{b,t}(w_t)-d_t^2A_{d_t}\}.
\]
The preceding finite identity gives, for every \(t\),
\[
\frac{
\sup_{c\in\mathcal M_t(G_t)}
\mathbb E_Z[(\widehat\tau_{w_t,t}-\tau_{n_t})^2]
}{
B^2A_{d_t}/m_t
}
=1+Y_t.
\]
Therefore convergence of the risk ratio to \(1\) is equivalent to convergence of \(Y_t\) to \(0\), by subtracting and adding the constant sequence \(1\). % lean: riskRatio_tendsto_iff_normalizedBlockExcess

\item Let
\[
X_t:=
\frac{1}{n_tA_{d_t}}
\sum_{i\in V_{n_t}}
\mathbb E_Z\!\left[
\{w_{i,t}(Z_{b(i)})-g_{d_t}(Z_{b(i)})\}^2
\right].
\]
The block-excess sum is nonnegative by the lower bound above.  The blockwise distance inequality gives
\[
Y_t\le 2\sqrt{X_t}+X_t,
\]
after using \(n_t=m_td_t\) and simplifying
\[
\frac{
2\sqrt{d_t^2A_{d_t}}\sqrt{m_t}\sqrt{d_tE_t}
+d_tE_t
}{
m_td_t^2A_{d_t}
}
=
2\sqrt{\frac{E_t}{n_tA_{d_t}}}
+\frac{E_t}{n_tA_{d_t}}.
\]
Thus \(X_t\to0\) implies \(Y_t\to0\) by squeezing, and the identity in the previous step gives the asserted convergence of the risk ratio. % lean: normalizedWeightDistance_tendsto_implies_riskRatio

\item Assume \(d_t\to\infty\).  For all sufficiently large \(t\), \(\beta<d_t\).  Define \(w_t\) by adding to the canonical weights a centered full-block perturbation on one distinguished unit in each block:
\[
w_{i,t}(Z)
=
g_{d_t}(Z_{b(i)})
+
\mathbf 1\{i=i_b^\star\}
\left\{
\frac{2d_tA_{d_t}}{(p(1-p))^{d_t}}
\right\}^{1/2}
\prod_{j\in V_{b(i)}}(Z_j-p).
\]
The full-block centered monomial has order \(d_t\), so it is orthogonal under \cref{obj:ass:bernoulli-design} to every monomial of order at most \(\beta\); consequently these weights remain in \(\mathcal E_t^{\mathrm{loc,lin}}\).  The exact identity above gives the lower bound \(1\) for the risk ratio.  The witness block bound gives, eventually,
\[
\Psi_{b,t}(w_t)
\le
d_t^2A_{d_t}
+2\sqrt{d_t^2A_{d_t}}\sqrt{2d_tA_{d_t}}
+2d_tA_{d_t},
\]
and therefore
\[
\frac{
\sup_{c\in\mathcal M_t(G_t)}
\mathbb E_Z[(\widehat\tau_{w_t,t}-\tau_{n_t})^2]
}{
B^2A_{d_t}/m_t
}
\le
1+2\sqrt{\frac{2}{d_t}}+\frac{2}{d_t}.
\]
Since \(d_t\to\infty\), squeezing proves the displayed risk-ratio convergence for this sequence. % lean: orthogonalWitness_riskRatio_tendsto_one

\item For the same sufficiently large \(t\), the perturbation has squared energy \(2d_tA_{d_t}\) on the distinguished unit of each block and zero perturbation energy on the other units.  Hence
\[
\sum_{i\in V_{n_t}}
\mathbb E_Z\!\left[
\{w_{i,t}(Z_{b(i)})-g_{d_t}(Z_{b(i)})\}^2
\right]
=
m_t(2d_tA_{d_t}),
\]
and division by \(n_tA_{d_t}=m_td_tA_{d_t}\) gives the value \(2\).  If \(\pi\) preserves block membership, the complete-block score is invariant under the induced within-block relabeling, and the full-block centered perturbation has the same squared energy after the relabeling.  The same calculation gives the relabeled normalized distance \(2\). % lean: orthogonalWitness_normalized_relabel_distance_eq_two
\end{enumerate}
\end{proof}
